\documentclass[11pt,a4paper]{article}
\usepackage{amsmath,amssymb,amsthm}
\usepackage{geometry}
\geometry{margin=2.5cm}
\usepackage{hyperref}
\usepackage{enumitem}

\newtheorem{theorem}{Theorem}[section]
\newtheorem{lemma}[theorem]{Lemma}
\newtheorem{proposition}[theorem]{Proposition}
\newtheorem{corollary}[theorem]{Corollary}
\newtheorem{remark}[theorem]{Remark}

\newcommand{\C}{\mathbb{C}}
\newcommand{\R}{\mathbb{R}}
\newcommand{\N}{\mathbb{N}}
\newcommand{\LOp}{\mathcal{L}}
\newcommand{\TOp}{\mathcal{T}}
\newcommand{\spec}{\operatorname{spec}}
\newcommand{\supp}{\operatorname{supp}}
\newcommand{\tr}{\operatorname{tr}}
\newcommand{\Tr}{\operatorname{Tr}}
\newcommand{\RE}{\operatorname{Re}}
\newcommand{\IM}{\operatorname{Im}}
\DeclareMathOperator{\detreg}{det_{(2)}}

\title{\bf Step 3: Spectral Representation of the Trace\\
(Regularised Stieltjes Transform of \(\mu_\sigma\))}
\author{}
\date{\today}

\begin{document}
\maketitle

\noindent
\textbf{Recap of Step~2.}  From the similarity theory of the GDST programme and the determinant identity, one obtains, after a careful algebraic manipulation, the following relation for the regularised resolvent trace of the self‑adjoint operator \(H_\sigma\):
\begin{equation}\label{eq:step2}
\Tr\!\Big[ (s - H_\sigma)^{-1} - s^{-1}I \Big] \;=\; \frac{\zeta'(s)}{\zeta(s)} + \Phi(s),
\qquad \Phi \text{ entire},
\end{equation}
valid for \(\RE s\) outside a neighbourhood of the spectrum of \(H_\sigma\).
(For the derivation see Step~2; the operator \(H_\sigma\) is unitarily equivalent, up to a finite‑rank perturbation, to the correlation operator \(T_\sigma\).)
The purpose of Step~3 is to bring the left‑hand side into the form of a regularised Stieltjes transform of a purely discrete positive measure, which will later permit the identification of the poles of \(\zeta'/\zeta\) with the support of this measure.

\section{Spectral decomposition of \(H_\sigma\)}

The operator \(H_\sigma\) (\(\sigma > 0\)) is a compact, positive, self‑adjoint trace‑class operator on the separable Hilbert space \(L^2[0,\pi]\).  Hence it possesses an orthonormal basis of eigenvectors \(\{e_i(\sigma)\}_{i=1}^\infty\) and a non‑increasing sequence of eigenvalues \(\{\lambda_i(\sigma)\}_{i=1}^\infty \subset \R_+\) satisfying
\[
\lambda_1(\sigma) \ge \lambda_2(\sigma) \ge \dots \ge 0, \qquad 
\lim_{i\to\infty} \lambda_i(\sigma) = 0, \qquad 
\sum_{i=1}^\infty \lambda_i(\sigma) < \infty .
\]
For any complex number \(z \notin \overline{\{\lambda_i(\sigma)\}}\) the resolvent of \(H_\sigma\) expands as
\[
(z - H_\sigma)^{-1} = \sum_{i=1}^\infty \frac{1}{z - \lambda_i(\sigma)} \, |e_i(\sigma)\rangle\langle e_i(\sigma)|,
\]
the series converging in operator norm and uniformly on compact sets that avoid the spectrum (see \cite{Simon}).

\section{Regularised resolvent trace and regularised Stieltjes transform}

Because the eigenvalues accumulate at zero, the series \(\sum_i 1/(s - \lambda_i(\sigma))\) diverges for generic \(s\).  However, the regularised resolvent
\[
R_{\mathrm{reg}}(s) \;:=\; (s - H_\sigma)^{-1} - s^{-1}I
\]
is trace‑class.  Indeed,
\[
R_{\mathrm{reg}}(s) = s^{-1} H_\sigma (s - H_\sigma)^{-1},
\]
and \(H_\sigma (s - H_\sigma)^{-1}\) is trace‑class (the trace‑class operators form an ideal).  Its trace is therefore well defined and given by the absolutely convergent series
\[
\Tr\big[ R_{\mathrm{reg}}(s) \big] \;=\; \sum_{i=1}^\infty \Bigl( \frac{1}{s - \lambda_i(\sigma)} - \frac{1}{s} \Bigr).
\]
Convergence follows from the estimate
\[
\Bigl| \frac{1}{s - \lambda_i} - \frac{1}{s} \Bigr|
= \frac{\lambda_i}{|s|\,|s - \lambda_i|}
\le \frac{\lambda_i}{|s|(|s| - \lambda_1)} \qquad (|s| > \lambda_1),
\]
and the summability of \(\sum \lambda_i\).

Define the \textbf{spectral measure} \(\mu_\sigma\) of \(H_\sigma\) as the positive discrete Borel measure
\[
\mu_\sigma \;:=\; \sum_{i=1}^\infty \delta_{\lambda_i(\sigma)},
\]
where \(\delta_x\) denotes the Dirac point mass at \(x\).  Then the regularised trace is precisely the \textbf{regularised Stieltjes transform} of \(\mu_\sigma\):
\begin{equation}\label{eq:Stieltjes}
\Tr\big[ R_{\mathrm{reg}}(s) \big] \;=\; \int_{\mathbb{R}}
\Bigl( \frac{1}{s - \lambda} - \frac{1}{s} \Bigr) \, d\mu_\sigma(\lambda)
\;=:\; S_{\mathrm{reg}}(s;\mu_\sigma).
\end{equation}
For \(|s| > \lambda_1(\sigma)\) the integral converges absolutely and defines a holomorphic function of \(s\).

\section{Conclusion of Step~3}

Inserting the expression \eqref{eq:Stieltjes} into the identity \eqref{eq:step2} yields the fundamental relation
\begin{equation}\label{eq:main3}
\boxed{\; S_{\mathrm{reg}}(s;\mu_\sigma) \;=\; \frac{\zeta'(s)}{\zeta(s)} + \Phi(s), \qquad \Phi \text{ entire}. \;}
\end{equation}

Both sides are meromorphic functions on the whole complex plane.  The left‑hand side has simple poles exactly at the points of \(\supp\mu_\sigma = \{\lambda_i(\sigma)\}\) (the regularisation removes the artificial singularity at \(s=0\)).  The right‑hand side has simple poles at the non‑trivial zeros of \(\zeta(s)\) (and at the pole \(s=1\), which can be treated separately).  Consequently, the sets of poles must coincide, a conclusion that will be made rigorous in Step~4.

\begin{remark}
The regularised Stieltjes transform \(S_{\mathrm{reg}}(s;\mu_\sigma)\) depends on the real parameter \(\sigma\) through the eigenvalues \(\lambda_i(\sigma)\).  Later (Step~5) we will show that the closed support of \(\mu_\sigma\) is in fact independent of \(\sigma\), a fact that forces the zeros of \(\zeta(s)\) to lie on the critical line \(\RE s = \frac12\).
\end{remark}

\begin{thebibliography}{1}
\bibitem{Simon} B.~Simon, \emph{Trace Ideals and Their Applications}, 2nd ed., AMS, 2005.
\end{thebibliography}

\end{document}